\documentclass[11pt,letterpaper]{article}
\usepackage[T1]{fontenc}
\usepackage{lmodern}
\usepackage[margin=1in]{geometry}
\usepackage{amsmath,amssymb,amsthm,mathtools}
\usepackage{microtype,booktabs,enumitem}
\usepackage{xurl}
\usepackage{xcolor}
\usepackage[colorlinks=true,linkcolor=blue!50!black,citecolor=blue!50!black,urlcolor=blue!50!black]{hyperref}
\hypersetup{pdftitle={Reciprocal Differential Energy on Finite Double Covers},pdfauthor={Drafted for Edward Baker}}
\numberwithin{equation}{section}
\newtheorem{theorem}{Theorem}[section]
\newtheorem{proposition}[theorem]{Proposition}
\newtheorem{lemma}[theorem]{Lemma}
\newtheorem{corollary}[theorem]{Corollary}
\theoremstyle{definition}
\newtheorem{definition}[theorem]{Definition}
\theoremstyle{remark}
\newtheorem{remark}[theorem]{Remark}
\newcommand{\C}{\mathbb C}
\newcommand{\R}{\mathbb R}
\newcommand{\D}{\mathbb D}
\newcommand{\E}{\mathcal E}
\newcommand{\Eo}{\mathcal E_0}
\newcommand{\dd}{\,d}
\newcommand{\ii}{\mathrm i}
\newcommand{\re}{\operatorname{Re}}
\newcommand{\im}{\operatorname{Im}}
\newcommand{\arsinh}{\operatorname{arsinh}}
\newcommand{\arcosh}{\operatorname{arcosh}}
\newcommand{\supp}{\operatorname{supp}}
\newcommand{\osc}{\mathrm{osc}}
\newcommand{\ext}{\mathrm{ext}}
\newcommand{\cut}{\mathrm{cut}}
\newcommand{\phys}{\mathrm{phys}}
\newcommand{\repo}{\texttt{shifted-zeta-positivity}}
\setlist[enumerate]{itemsep=3pt,topsep=5pt}
\setlength{\emergencystretch}{2em}
\allowdisplaybreaks[2]
\title{Reciprocal Differential Energy on Finite Double Covers\\
\large Boundary reductions, quadratic models, and the theta heat flow}
\author{Drafted for Edward Baker}
\date{Draft manuscript --- 22 October 2026}
\begin{document}
\maketitle
\begin{abstract}
For a holomorphic function on a bounded spatial domain, the differential
$dz/\sqrt H$ on its two-sheeted cover has energy $2\int |H|^{-1}\dd A$.
This energy is finite at a simple zero and diverges at a multiple zero.
The research objective is an independent quantitative upper estimate for
the canonical theta heat flow, rather than another formulation of
simplicity. We develop the finite-domain analytic framework for this
objective. On a two-branch annulus, a single-valued real primitive gives
an exact reduction to a projected relative period and a complete exterior
boundary correction. Subtracting a quadratic differential with the same
core period yields explicit energy, relative-period, and full-boundary
remainders bounded independently of the branch separation. A complete
moment-controlled chart at a hypothetical nondegenerate heat event supplies
all constants and a finite upper bound; a sharper bound strictly improves
the initial coarse certificate on an explicit window. An exact quadratic
calculation shows why a bounded core period and bounded exterior norm
budgets do not bound the full energy. A previously recorded interval
computation on one regular theta fiber supports the finite acquisition
framework. The analytic estimates remain compatible with ordinary double
collisions: no endpoint-uniform theta estimate, RH theorem, or simplicity
exclusion is obtained.
\end{abstract}
\tableofcontents
\section{Project goals and the results established so far}\label{sec:goals}

The project studies the reciprocal differential on finite covers of the
canonical theta heat field. The intended long-term application is to
control degeneration as zeros approach one another. The distinction between
an identity and an estimate is central. Expressing simplicity as local
integrability does not prove simplicity. A useful advance would be an
upper estimate obtained from the unchanged canonical source and remaining
controlled in the relevant limiting regime, or a demonstrable finite
acquisition advantage over direct integration using the same information.

This manuscript develops the analytic part of that program. Its setting
is a bounded disk containing two simple real zeros. The resulting cover
is an annulus, with two physical boundary components, one absolute core
cycle, and a relative path connecting the two boundary components. The
relative period must be retained: it can diverge even when the core period
is constant and the exterior differential has uniformly bounded traces.

\subsection{Canonical normalization and the motivating question}
We use
\begin{align}
 H(t,z)&=\int_0^\infty e^{tu^2}\Phi(u)\cos(zu)\dd u,\label{eq:Hdef}\\
 \Phi(u)&=\sum_{n=1}^\infty
 (2\pi^2n^4e^{9u}-3\pi n^2e^{5u})e^{-\pi n^2e^{4u}}.
 \label{eq:Phi}
\end{align}
The classical theta identity fixes the normalization
\begin{equation}\label{eq:xi}
 H(0,z)=\frac18\xi\left(\frac12+\frac{\ii z}{2}\right),\qquad
 \xi(s)=\frac{s(s-1)}2\pi^{-s/2}\Gamma(s/2)\zeta(s).
\end{equation}
Appendix~\ref{sec:normalization} derives this identity and the complete
source reflection $\Phi(u)=\Phi(-u)$ from the Gaussian theta identity.
The convention agrees with \cite[(1)--(4)]{RT}.
Our derivative and energy estimates below start from
\eqref{eq:Hdef}--\eqref{eq:Phi}. In particular,
\begin{equation}\label{eq:heat}
 \partial_t H=-\partial_z^2H.
\end{equation}
The RH motivation concerns the zero set of \eqref{eq:xi}; none of the local
theorems below assumes RH. In a hypothetical ordinary collision, the
zeros split at the square-root scale in positive physical time. The
reciprocal energy then grows logarithmically. Preventing such a collision
would require a source estimate incompatible with that growth, with
constants valid in the limiting family. The present finite estimates do
not have that property.

\subsection{Analytic contributions}
The results are organized as follows.
\begin{enumerate}
\item Section~\ref{sec:geometry} constructs the cover, normalizes the
energy, and proves the local integrability statement at all multiplicities.
It gives explicit coordinates for the full bordered annulus.
\item Section~\ref{sec:boundary} proves the complete cut-boundary identity
and a real-primitive reduction. With consistently oriented periods $q,p$,
\begin{equation}\label{eq:roadmapGreen}
 \E=\im(\bar q p)+R_{\ext},\qquad
 |R_{\ext}|\le 2\pi^2K_\partial^2.
\end{equation}
The bound uses the exterior differential trace; the relative period
contains the missing degeneration.
\item Section~\ref{sec:quadratic} evaluates the entire quadratic cover and
each of its cut and physical-boundary contributions. It proves a scoped
failure of upper estimates based only on fixed exterior norm budgets and
the permitted core period.
\item Section~\ref{sec:remainders} subtracts a quadratic differential with
the same core period. It proves explicit bounds for the full energy,
relative period, and both boundary contributions, with constants independent
of the root separation when the unit and its derivatives are controlled.
\item Section~\ref{sec:heat} derives a complete quantitative chart from
canonical source moments at a hypothetical ordinary event. It proves a
finite logarithmic upper bound and a strict improvement over its coarse
version on a specified smaller window. These are shared direct estimates,
not a new restriction on the source's admissible events.
\end{enumerate}
Section~\ref{sec:enclosures} records analytic error consequences. Numerical
support is confined to Section~\ref{sec:support}, with the detailed
computational exposition left to its cited record. Section~\ref{sec:open}
identifies the remaining endpoint and coverage problems.

\subsection{Comparison principle and scope}
The direct comparator is always the same quantity,
$2\int_\Omega |H|^{-1}\dd A$, supplied with all the same source,
normalization, root, boundary, and period observations. It may use the
same coordinate substitutions and exact identities. Consequently, a
better bound than one coarse direct certificate does not establish a
general advantage over direct integration. Likewise, a native numerical
implementation can be faster without having more information.

The results collected here consolidate \cite{Notes}; their proofs are
included in full. Classical complex analysis, the exact canonical
normalization above, and ordinary calculus supply the background. There
is no compact finite-genus model of the complete entire theta cover in
the argument, and no omission of the spatial boundary of a finite cover.

\section{Finite double covers and reciprocal energy}\label{sec:geometry}

\subsection{Regular fibers and the factor of two}
Let $\Omega\subset\C$ be a bounded domain with smooth boundary, and let
$H$ be holomorphic on a neighborhood of $\overline\Omega$, nonzero on
$\partial\Omega$. When every zero in $\Omega$ is simple, set
\begin{equation}\label{eq:cover}
 X=\{(z,w):z\in\overline\Omega,\ w^2=H(z)\},\qquad
 \omega=\frac{dz}{w}.
\end{equation}
The complex structure on $X$ gives its orientation. At a zero $\alpha$,
write $H(z)=(z-\alpha)U(z)$, $U(\alpha)\ne0$. In the local coordinate
$z=\alpha+\zeta^2$,
\begin{equation}\label{eq:branchregular}
 w=\zeta\sqrt{U(\alpha+\zeta^2)},\qquad
 \omega=\frac{2\dd\zeta}{\sqrt{U(\alpha+\zeta^2)}}.
\end{equation}
Thus the branch point is regular on the cover and $\omega$ is holomorphic.

\begin{proposition}[Full energy]\label{prop:energy}
The conformally invariant energy of $\omega$ is
\begin{equation}\label{eq:energy}
 \E(H;\Omega):=\frac{\ii}{2}\int_X\omega\wedge\bar\omega
       =2\int_\Omega\frac{\dd A_z}{|H(z)|}.
\end{equation}
The factor two counts both sheets. In a local conformal coordinate,
the energy convention has no additional factor $1/2$.
\end{proposition}
\begin{proof}
Away from the finitely many branch points, $z$ is a coordinate on either
sheet and $\omega=f(z)\dd z$ with $|f|^2=|H|^{-1}$.
Since $(\ii/2)\dd z\wedge\dd\bar z=\dd A_z$, each sheet contributes
$\int_\Omega |H|^{-1}\dd A$. The branch points have area zero, and
\eqref{eq:branchregular} gives local finiteness, so the equality extends
across them.
\end{proof}

\begin{proposition}[Multiplicity and local integrability]\label{prop:multi}
If $H(z)=(z-z_0)^kU(z)$ with $U(z_0)\ne0$, its local reciprocal energy
is finite precisely when $k=1$ (or when $k=0$, meaning no zero).
A double zero gives logarithmic divergence; a zero of order $k>2$
gives a power divergence.
\end{proposition}
\begin{proof}
On a sufficiently small disk there are constants $0<u_-\le|U|\le u_+$.
The radial energy is comparable to
\begin{equation}\label{eq:radial}
 \int_0^\varepsilon r^{1-k}\dd r.
\end{equation}
This converges for $k<2$, is logarithmic for $k=2$, and diverges more
strongly for $k>2$.
The same conclusion holds on the normalization of the singular cover.
For $k=2\ell$, its two normalized branches have
$w=\pm(z-z_0)^\ell\sqrt U$, and $\omega$ has a pole of order $\ell$.
For $k=2\ell+1$, put $z-z_0=\zeta^2$; then
$\omega=2\zeta^{-2\ell}U(z_0+\zeta^2)^{-1/2}\dd\zeta$.
It is regular for $k=1$ and has a pole for each odd $k\ge3$.
These expressions also explain geometrically why the singular fiber
cannot simply be included among the smooth finite-energy fibers.
\end{proof}

For a nonzero holomorphic field on a closed bounded domain with no boundary
zero, local reciprocal integrability therefore restates simplicity of its
interior zeros. We use this observation once, as a classification of the
observable; it is not an independent simplicity estimate.

\subsection{A disk with two real branches}
In the remainder of the finite-cover analysis let
\begin{equation}\label{eq:diskpair}
 \Omega=D(C,R):=\{|z-C|<R\},\quad C\in\R,\quad
 \alpha<\beta\in\Omega\cap\R,
\end{equation}
and suppose these are exactly the two zeros of $H$, both simple. Then
\begin{equation}\label{eq:unitfactor}
 H(z)=(z-\alpha)(z-\beta)V(z),\quad
 m=\frac{\alpha+\beta}{2},\quad a=\frac{\beta-\alpha}{2}>0.
\end{equation}
The unit $V$ is holomorphic and nonvanishing on the disk and its boundary.
Choose a holomorphic square root of $V$ and put
\begin{equation}\label{eq:eta}
 y_0^2=(z-m)^2-a^2,\quad b=V^{-1/2},\quad
 \eta=\frac{dz}{y_0},\quad \omega=b(z)\eta.
\end{equation}
The two covers are identified by $w=\sqrt V\,y_0$.

\begin{lemma}[Explicit annulus coordinates]\label{lem:cylinder}
The full cover is an annulus. It is parametrized by
\begin{equation}\label{eq:Joukowski}
 z=m+a\cosh u,\qquad y_0=a\sinh u,\qquad
 u=s+\ii\theta\pmod{2\pi\ii},\qquad \eta=\dd u.
\end{equation}
Writing $\delta=m-C$, its domain is $-U(\theta)\le s\le U(\theta)$,
where
\begin{equation}\label{eq:Ushift}
 U(\theta)=\arcosh\left(
 \frac{-\delta\cos\theta+
 \sqrt{R^2+(a^2-\delta^2)\sin^2\theta}}a\right)>0.
\end{equation}
Both physical boundary components are retained.
\end{lemma}
\begin{proof}
In the nonzero coordinate $\zeta=e^u$,
$z=m+(a/2)(\zeta+\zeta^{-1})$ and
$y_0=(a/2)(\zeta-\zeta^{-1})$. The inverse is
$\zeta=(z-m+y_0)/a$, nonzero because
$(z-m+y_0)(z-m-y_0)=a^2$. Thus these are genuine coordinates on
the cover, including $u=0,\pi\ii$, where $y_0=0$ but the derivative
of the pair $(z,y_0)$ is nonzero.

The real interior roots imply $|\delta|+a<R$. Set $v=\cosh s\ge1$.
A direct calculation gives
\[
 |z-C|^2=a^2v^2+2a\delta v\cos\theta+
                \delta^2-a^2\sin^2\theta.
\]
The equation $|z-C|=R$ has the positive solution for $v$ in
\eqref{eq:Ushift}. At $s=0$, $z=m+a\cos\theta$ lies strictly inside
the disk. The quadratic expression has just one crossing for $v\ge1$,
so the domain is precisely the asserted strip. Its positive continuous
width and $2\pi$ periodicity identify an annulus, with boundaries
$s=\pm U(\theta)$. Differentiating \eqref{eq:Joukowski} gives
$dz/y_0=\dd u$, with the identity extending at the branch points.
\end{proof}

We label the boundary component over $z=C+R$ with $y_0>0$ as the plus
component. Its induced positive orientation projects counterclockwise
onto the base circle; its core period for $\eta$ is $2\pi\ii$.
The other component has the opposite sheet label, but its induced
orientation also projects counterclockwise. The sheet involution
$\iota$ satisfies $\iota^*\eta=-\eta$ and $\iota^*\omega=-\omega$.

\begin{corollary}[Full direct coordinate formula]\label{cor:direct}
If $H(\bar z)=\overline{H(z)}$, the energy can be written without a
branch singularity as
\begin{equation}\label{eq:directU}
 \E=4\int_0^\pi\int_0^1
 \frac{U(\theta)}{|V(m+a\cosh(tU(\theta)+\ii\theta))|}
          \dd t\dd\theta.
\end{equation}
\end{corollary}
\begin{proof}
In \eqref{eq:Joukowski}, the energy density is
$|b(z)|^2\dd s\dd\theta=|V(z)|^{-1}\dd s\dd\theta$.
The replacements $s\mapsto-s$ and $\theta\mapsto2\pi-\theta$
conjugate $z$, while leaving $U$ unchanged in the needed pairs.
Reality of $V$ gives the two factors of two. Finally put $s=tU(\theta)$;
the area Jacobian in $(s,\theta)$ is exactly $U(\theta)$.
\end{proof}

\section{Complete boundary identities}\label{sec:boundary}

\subsection{The complex primitive and paired cuts}
Cut the annulus along a simple relative arc connecting its two physical
boundary components, obtaining a disk $D$. Choose a single primitive
$F$ on this disk, $\dd F=\omega$, with one common additive constant.
\begin{proposition}[Complete Stokes identity]\label{prop:stokes}
With all cut edges and both physical boundaries given their induced
orientations,
\begin{equation}\label{eq:stokes}
 \E=\frac{\ii}{2}\int_{\partial D}F\bar\omega.
\end{equation}
For a paired cut edge, let $e_+$ have its $\partial D$ orientation and
identify $e_-$ oppositely under gluing. Then
\begin{equation}\label{eq:jump}
 \int_{e_+\cup e_-}F\bar\omega
       =J_e\overline{P_e},\qquad
 J_e=F_+-F_-,\quad P_e=\int_{e_+}\omega.
\end{equation}
The jump $J_e$ is constant; every unpaired physical arc remains in
\eqref{eq:stokes}.
\end{proposition}
\begin{proof}
Holomorphicity implies $\dd\bar\omega=0$, so
$\dd(F\bar\omega)=\omega\wedge\bar\omega$. Stokes proves
\eqref{eq:stokes}. On a paired edge the differential agrees under gluing
and the orientations are opposite. Its two integrals therefore combine
to $\int_{e_+}(F_+-F_-)\bar\omega$. The difference has zero differential,
which proves both constancy and \eqref{eq:jump}.
\end{proof}

On a compact boundaryless surface, a symplectic polygon pairs all edges;
the contributions are $A_j\bar B_j-B_j\bar A_j$, and hence
$\E=\sum_j\im(\bar A_jB_j)$. This is the usual bilinear relation
\cite[Section 3.1.2]{Grava}. On the present annulus there are unpaired
physical boundaries. It has a relative transverse path, not a compact
genus-one dual cycle. Capping those boundaries without extending the
differential and accounting for its poles would not justify deleting them.

\subsection{A real primitive on the whole annulus}
Let $q=\int_c\omega$ on the counterclockwise plus core. Let $p$ be the
integral along a fixed relative path from the plus basepoint over
$z_*=C+R$ to its minus partner. Choose the path and core consistently;
in the real two-root chart we take the path through $\beta$.
Changing the relative path by a core adds an integer multiple of $q$,
which leaves $\im(\bar q p)$ unchanged.

\begin{theorem}[Real-primitive reduction]\label{thm:realprimitive}
Suppose $q\ne0$ and define
\begin{equation}\label{eq:lambdaS}
 \lambda=\frac{\ii\bar q}{|q|},\quad
 P=\re(\lambda p),\quad S=\im(\bar q p)=-P|q|.
\end{equation}
On the plus boundary, parametrized by $z=C+Re^{\ii\theta}$, write
$\omega=B(\theta)\dd\theta$ and set
\begin{equation}\label{eq:vt}
 v(\theta)=\re\left(\lambda\int_0^\theta B(s)\dd s\right),
 \qquad t(\theta)=\im(\lambda B(\theta)).
\end{equation}
Then
\begin{equation}\label{eq:realGreen}
 \boxed{\E=S+R_{\ext},\qquad
 R_{\ext}=2\int_0^{2\pi}v(\theta)t(\theta)\dd\theta.}
\end{equation}
If $|B|\le K_\partial$ on the whole circle, then
\begin{equation}\label{eq:Rbound}
 |R_{\ext}|\le2\pi^2K_\partial^2.
\end{equation}
If $q=0$, choose any unit $\lambda$; the same formula holds with
$S=0$ and with \eqref{eq:vt} defined using that choice.
\end{theorem}
\begin{proof}
Every loop in the annulus is homologous to an integer multiple of $c$.
Since $\lambda q=\ii|q|$, the real part of every period of
$\lambda\omega$ is zero. Thus
$h=\re(\lambda F)$ is a single-valued harmonic function on $X$.
Its differential changes sign under $\iota$, so $h\circ\iota+h$ is
constant. A single global change of additive constant makes $h$ odd.
At the two basepoints,
$h_--h_+=\re(\lambda p)=P$ and $h_-=-h_+$, giving $h_+=-P/2$.
Consequently the boundary traces are $h_+=-P/2+v$ and $h_-=-h_+$.

For completeness, in a local coordinate $x+\ii y$ write
$\lambda\omega=f\dd z$ and a primitive $h+\ii k$. The Cauchy--Riemann
equations give $h_x=\re f$, $h_y=-\im f$. Hence
$|\nabla h|^2\dd A=(\ii/2)\omega\wedge\bar\omega$.
On a positively oriented boundary the outward normal flux of $h$
is $\dd k=\im(\lambda\omega)$. Green's formula therefore gives
\[
 \E=\int_X|\nabla h|^2\dd A
       =\int_{\partial X}h\,\im(\lambda\omega).
\]
Both sheet values change sign under $\iota$, so the two physical
components contribute equally. Using $\int_0^{2\pi}t=|q|$ yields
\[
 \E=2\int_0^{2\pi}(-P/2+v)t\dd\theta
      =-P|q|+2\int_0^{2\pi}vt\dd\theta.
\]
This proves \eqref{eq:realGreen} and its sign convention.

The function $v$ is periodic, with $v(0)=v(2\pi)=0$ and
$|v'|\le K_\partial$. Integrating from either endpoint gives
$|v(\theta)|\le K_\partial\min(\theta,2\pi-\theta)$.
Since the integral of this minimum is $\pi^2$, \eqref{eq:Rbound} follows.
If $q=0$, the complex primitive itself is single-valued and the preceding
argument applies with $\int t=0$, proving the stated extension.
\end{proof}

\begin{remark}[What disappeared from the cut formula]
The cut terms vanish for the real primitive because $h$ is single-valued.
They have not become zero information: the relative constant $P$ remains
in $S$. In particular, independently resetting the primitive constant on
each physical boundary would destroy \eqref{eq:realGreen}.
\end{remark}

\begin{proposition}[Real allocation between cuts and physical boundary]
\label{prop:cutallocation}
Choose the relative cut as above and the consistent polygon primitive.
The real cut contribution to \eqref{eq:stokes} is $S/2$, and the real
physical-boundary contribution is $S/2+R_{\ext}$.
\end{proposition}
\begin{proof}
Orient one copy of the cut from minus to plus, opposite to the orientation
of $p$. With the positive core chosen above, its primitive jump is $-q$.
Equation~\eqref{eq:jump} gives cut integral $q\bar p$, whose real energy
contribution is
$\re((\ii/2)q\bar p)=\im(\bar q p)/2=S/2$.
Subtracting this contribution from \eqref{eq:realGreen} proves the second
claim. An individual complex subtotal need not be real; only its real
energy contribution is asserted here.
\end{proof}

\begin{corollary}[An exterior floor and a finite energy interval]
If $|H|\ge\mu>0$ on the spatial circle, then
$K_\partial\le R/\sqrt\mu$, and
\begin{equation}\label{eq:periodInterval}
 \E\in\left[S-\frac{2\pi^2R^2}{\mu},
             S+\frac{2\pi^2R^2}{\mu}\right]\cap[0,\infty).
\end{equation}
\end{corollary}
\begin{proof}
The pullback satisfies $|B|=R/\sqrt{|H|}$; apply the theorem.
\end{proof}

If the field has real type, the unit is positive on the real interval,
and its square root is chosen positive there, then $q=\ii Q$ with
$Q>0$ (proved in Section~\ref{sec:remainders}). Thus $\lambda=1$.
Conjugation of the boundary trace gives
$B(2\pi-\theta)=-\overline{B(\theta)}$ and
$v(2\pi-\theta)=v(\theta)$, so
\begin{equation}\label{eq:halfboundary}
 Q=2\int_0^\pi\im B\dd\theta,\qquad
 R_{\ext}=4\int_0^\pi v\im B\dd\theta.
\end{equation}
These are exact symmetry reductions, not separate estimates of positive
sectors.

\section{The complete quadratic calibration}\label{sec:quadratic}

\subsection{Exact energy and both boundary payments}
Take $H_a(z)=z^2-a^2$ on $|z|<R$, with $0<a<R$. In
Lemma~\ref{lem:cylinder}, $m=C=0$ and
\begin{equation}\label{eq:Uround}
 U_a(\theta)=\arsinh\sqrt{(R/a)^2-\cos^2\theta},\qquad
 U_0:=U_a(0)=\arcosh(R/a).
\end{equation}
Here $\omega=\dd u$ and the full energy is the cylinder area.
\begin{proposition}[Exact quadratic calculation]\label{prop:quadratic}
One has
\begin{align}
 \E_a&=2\int_0^{2\pi}U_a(\theta)\dd\theta,
 \label{eq:Equad}\\
 4\pi\arcosh(R/a)&\le\E_a\le4\pi\arsinh(R/a),
 \label{eq:quadbounds}\\
 \E_a&=4\pi\log(2R/a)+O(a^2/R^2)\quad(a/R\longrightarrow0).
 \label{eq:quadasymp}
\end{align}
With $q=2\pi\ii$ and $p=-2U_0$ oriented plus to minus, the real
cut and physical contributions are
\begin{equation}\label{eq:quadpayments}
 E_{\cut}=2\pi U_0,\qquad
 E_{\phys}=\E_a-2\pi U_0.
\end{equation}
Each contains half of the leading logarithmic divergence.
\end{proposition}
\begin{proof}
The coordinate $u=s+\ii\theta$ has energy form
$(\ii/2)\dd u\wedge\dd\bar u=\dd s\dd\theta$, proving
\eqref{eq:Equad}. The minimum and maximum of $U_a$ are
$\arcosh(R/a)$ and $\arsinh(R/a)$, proving \eqref{eq:quadbounds}.
The identities
$\arcosh x=\log(x+\sqrt{x^2-1})$ and
$\arsinh x=\log(x+\sqrt{x^2+1})$ give the stated asymptotics.

To verify every boundary contribution directly, cut the cylinder at
$\theta=0$ and use $F=u$ on $0\le\theta\le2\pi$.
The bottom runs from $-U_0$ to $U_0$, the right side upward,
the top in reverse, and the left side downward. The bottom-to-top jump
is $-2\pi\ii$, so the paired cut integral is $-4\pi\ii U_0$.
Its energy contribution is $2\pi U_0$.
On the right $u=U_a+\ii\theta$ and on the left
$u=-U_a+\ii\theta$, with opposite $\theta$ orientations. Thus
\begin{align*}
 \int_{\partial X}F\bar\omega
 &=2\ii\int_0^{2\pi}(\theta U_a'-U_a)\dd\theta\\
 &=4\pi\ii U_0-4\ii\int_0^{2\pi}U_a\dd\theta.
\end{align*}
Multiplication by $\ii/2$ proves \eqref{eq:quadpayments}. A common
global primitive shift changes neither grouped contribution, since the
integral of $\bar\omega$ over either grouped boundary pair is zero.
\end{proof}

For the heat control $H(\tau,z)=z^2-2\tau$, $a=\sqrt{2\tau}$, so
\begin{equation}\label{eq:quadheat}
 \E_\tau=2\pi\log(1/\tau)+O(1)
 \qquad(\tau\downarrow0).
\end{equation}
The finite core period is not a finite-energy endpoint statement.

\subsection{A scoped limitation of bounded exterior summaries}
\begin{proposition}\label{prop:compression}
Fix $R>0$ and $0<a\le R/2$. The core period, an exterior field floor,
every prescribed finite collection of exterior differential derivative
upper bounds, and separately normalized boundary primitive oscillation
upper bounds can all be chosen uniformly in $a$. Nevertheless $\E_a$
is unbounded as $a\downarrow0$. No finite upper estimate for the energy
can depend only on those fixed common budget inequalities.
\end{proposition}
\begin{proof}
The core period is identically $2\pi\ii$. On $|z|=R$,
$|H_a|\ge3R^2/4$. On a fixed collar, for example
$3R/4\le|z|\le5R/4$, the expression
\[
 \frac1{\sqrt{z^2-a^2}}=\frac1z
       \left(1-\frac{a^2}{z^2}\right)^{-1/2}
\]
is holomorphic with the branch near one and uniformly bounded.
Cauchy's estimates on smaller fixed collars bound every fixed finite
derivative collection. In the cylinder parametrization the separately
rebased boundary primitives have increments
$\pm(U_a(\theta)-U_0)+\ii\theta$; they are uniformly bounded since
\[
 0\le U_a-U_0\le\arsinh(R/a)-\arcosh(R/a).
\]
But \eqref{eq:quadbounds} diverges. The missing relative constant is
$2U_0$, which is unbounded.
\end{proof}

The scope is important. Exact attained norms or complete exact boundary
values can encode $a$ and permit its recovery. The proposition concerns
fixed common upper-budget inequalities, not every possible use of exterior
data. It also shows why an absolute field error set admitting arbitrarily
small positive $a$ admits arbitrarily large finite reciprocal energy,
even if each admitted positive-$a$ fiber is regular.

\subsection{A lower bound for general pair charts}
\begin{lemma}\label{lem:screenlower}
Suppose $H(z)=(z^2-a^2)V(z)$ on $|z|\le R$, $0<2a<R$, and
$|V|\le M$. Then
\begin{equation}\label{eq:lowerlog}
 \E\ge\frac{16\pi}{5M}\log\frac{R}{2a}.
\end{equation}
Consequently, an independently obtained upper $\E\le B$ implies
$a\ge(R/2)\exp[-5MB/(16\pi)]$.
\end{lemma}
\begin{proof}
On $2a\le|z|\le R$, $|z^2-a^2|\le(5/4)|z|^2$.
Integrating $2\cdot4/(5M|z|^2)$ over this annulus proves
\eqref{eq:lowerlog}; rearranging proves the consequence.
\end{proof}
An endpoint conclusion requires a fixed positive $R$ and uniform $M,B$.
The upper estimates later in this manuscript do not supply uniform $B$.

\subsection{Capacity as supporting analysis}
On the round annulus $r<|z|<R$, the function
$\varphi(z)=\log(|z|/r)/\log(R/r)$ has boundary values zero and one.
Direct polar integration gives, with no factor $1/2$,
\[
 \operatorname{cap}A(r,R)=\int|\nabla\varphi|^2\dd A
        =\frac{2\pi}{\log(R/r)}.
\]
For a closed real one-form of core period $p$, circlewise
Cauchy--Schwarz and radial integration give
\[
 \int|\alpha|^2\dd A\ge\frac{p^2}{2\pi}\log(R/r).
\]
Indeed its tangential integral is $p$ on each circle, so its energy on
that circle is at least $p^2/(2\pi s)$. These formulas explain the
degeneration scale but do not give an upper modulus bound. In particular,
a trial potential gives an upper capacity bound and hence a lower bound
for the fixed-period core energy, not the upper estimate sought here.

\section{Quadratic subtraction with a bounded full remainder}\label{sec:remainders}

Continue with the disk and unit factor in \eqref{eq:diskpair}--\eqref{eq:eta}.
Assume throughout this section that
\begin{equation}\label{eq:unitbudgets}
 |V|\ge\nu>0,\qquad |b'|\le L,\qquad |b''|\le K
 \quad\hbox{on }\overline{D(C,R)},\qquad b=V^{-1/2}.
\end{equation}
These are explicit budgets on a fixed disk. Uniformity below means
independence of the branch separation within a family for which the same
$R,\nu,L,K$ work. Section~\ref{sec:heat} derives such budgets from the
complete canonical source on an event-dependent chart.

\subsection{The mean core coefficient and an exact primitive}
\begin{lemma}[The core mean]\label{lem:coremean}
With the plus marking of Section~\ref{sec:geometry},
\begin{equation}\label{eq:cmean}
 c:=\frac{q}{2\pi\ii}
   =\frac1\pi\int_0^\pi b(m+a\cos\theta)\dd\theta.
\end{equation}
Writing $b_0=b(m)$, one has
\begin{equation}\label{eq:cerror}
 |c|,|b_0|\le\nu^{-1/2},\qquad |c-b_0|\le\frac{Ka^2}{4}.
\end{equation}
If $V$ is real and nonzero on the slit $[\alpha,\beta]$, then $c\ne0$.
\end{lemma}
\begin{proof}
Shrink the plus contour counterclockwise onto the slit. On its upper
bank $y_0=\ii\sqrt{a^2-(z-m)^2}$ and the direction is from $\beta$
to $\alpha$; the lower bank has the opposite square root and direction.
Their contributions add to
\[
 q=2\ii\int_\alpha^\beta
       \frac{b(x)}{\sqrt{a^2-(x-m)^2}}\dd x.
\]
The endpoint arcs vanish in the local regular coordinate
\eqref{eq:branchregular}. Substituting $x=m+a\cos\theta$ proves
\eqref{eq:cmean}. The bound for $c$ follows by averaging $|b|$.
Taylor's formula on the real segment gives
\[
 b(m+h)=b_0+h b'(m)+h^2\int_0^1(1-t)b''(m+th)\dd t.
\]
The mean of $\cos\theta$ is zero and that of $\cos^2\theta$ is $1/2$,
proving the second estimate. For real nonzero $V$, its sign is constant
on the slit. The coherent $b$ has one fixed phase times a strictly
positive real function there, so its average cannot vanish. For a general
complex unit no such nonvanishing assertion is made.
\end{proof}

\begin{proposition}[Exact quadratic subtraction]\label{prop:subtract}
There is a single-valued holomorphic function $G$ on the whole cover,
normalized by $G(\beta)=0$, such that
\begin{equation}\label{eq:Gdecomp}
 \omega=c\eta+\dd G,\qquad G=y_0 f(z),
\end{equation}
where $f$ is holomorphic on the base disk and satisfies
\begin{equation}\label{eq:fode}
 ((z-m)^2-a^2)f'(z)+(z-m)f(z)=b(z)-c.
\end{equation}
If $\Eo$ is the full energy of $\eta$, then
\begin{equation}\label{eq:exactsubtract}
 \E=|c|^2\Eo+
 2\re\left[\bar c\frac{\ii}{2}\int_{\partial X}G\bar\eta\right]
       +\frac{\ii}{2}\int_{\partial X}G\dd\bar G.
\end{equation}
Here $\partial X$ consists of the two physical boundaries.
\end{proposition}
\begin{proof}
The only period of $\omega-c\eta$ is $q-2\pi\ii c=0$. Its path
integral is therefore a single-valued holomorphic primitive $G$.
The differential changes sign under sheet interchange, so
$G\circ\iota+G$ is constant. Its value at the fixed point $\beta$
is zero, proving that $G$ is odd and, in particular, also vanishes at
$\alpha$. The quotient $G/y_0$ is even and descends away from the two
branches. At either branch use $z-\gamma=\zeta^2$. Both $G$ and
$y_0$ are odd in $\zeta$, while $y_0/\zeta$ is nonzero. Their quotient
extends holomorphically. Differentiation of $G=y_0 f$ proves
\eqref{eq:fode}.

Expand the energy of $c\eta+\dd G$. The two mixed integrals are complex
conjugates. For one of them,
$\dd(G\bar\eta)=\dd G\wedge\bar\eta$, and the last term is
$\dd(G\dd\bar G)=\dd G\wedge\dd\bar G$. Stokes gives
\eqref{eq:exactsubtract}. These correction forms are single-valued;
paired cut edges cancel, so no relative or physical contribution was
suppressed. The mixed term in \eqref{eq:exactsubtract} is essential.
\end{proof}

\subsection{Uniform energy estimates}
Set $\delta=m-C$, $R_-=R-|\delta|$, $R_+=R+|\delta|$. Disk inclusion
about $m$ and \eqref{eq:quadbounds} give
\begin{equation}\label{eq:shiftedsandwich}
 4\pi\arcosh(R_-/a)\le\Eo
 :=2\int_{D(C,R)}\frac{\dd A}{|z-\alpha||z-\beta|}
 \le4\pi\arsinh(R_+/a).
\end{equation}
The roots are interior, so $R_->a$ and $R_+<2R$.

\begin{theorem}[Bounds independent of the gap]\label{thm:remainder}
Under \eqref{eq:unitbudgets}, define
\begin{align}
 C_{\osc}&=\frac{8\pi RL}{\sqrt\nu}+8\pi R^2L^2,\label{eq:Cosc}\\
 C_E&=C_{\osc}+\frac{8\pi KR^2}{\sqrt\nu}.\label{eq:CE}
\end{align}
Then the complete energy satisfies
\begin{equation}\label{eq:energyremainders}
 |\E-|b_0|^2\Eo|\le C_{\osc},\qquad
 |\E-|c|^2\Eo|\le C_E.
\end{equation}
\end{theorem}
\begin{proof}
For any $\gamma$ in the disk, let $\rho(\theta)$ be the exit distance
along the ray from $\gamma$ in direction $\theta$. Polar integration
and Cauchy--Schwarz give
\[
 \int_{D(C,R)}\frac{\dd A}{|z-\gamma|}
 =\int_0^{2\pi}\rho(\theta)\dd\theta
 \le\sqrt{2\pi\int_0^{2\pi}\rho(\theta)^2\dd\theta}=2\pi R.
\]
Put $d_1=|z-\alpha|$, $d_2=|z-\beta|$. Since
$|z-m|\le(d_1+d_2)/2$ and $|z-m|<2R$,
\begin{equation}\label{eq:weightedI}
 I_1:=\int\frac{|z-m|}{d_1d_2}\dd A\le2\pi R,
 \qquad
 I_2:=\int\frac{|z-m|^2}{d_1d_2}\dd A\le4\pi R^2.
\end{equation}
All integrals here are over the full base disk. The straight segment
from $m$ to $z$ remains in the disk, so $|b(z)-b_0|\le L|z-m|$.
Expanding the square, including its mixed term, yields
\[
 \big||b(z)|^2-|b_0|^2\big|
 \le2\nu^{-1/2}L|z-m|+L^2|z-m|^2.
\]
Multiply by $2/(d_1d_2)$ and apply \eqref{eq:weightedI}; this proves
the first bound. By \eqref{eq:cerror},
\[
 \big||c|^2-|b_0|^2\big|\le\frac{Ka^2}{2\sqrt\nu}.
\]
Furthermore $\arsinh x\le x$ for $x\ge0$, whence
$a^2\arsinh(R_+/a)\le aR_+\le R_+^2\le4R^2$.
Using the upper bound in \eqref{eq:shiftedsandwich} proves the second
estimate with exactly the constant \eqref{eq:CE}.
\end{proof}

\subsection{The relative period and the full boundary remainder}
Let $\ell=C+R-m>a$. For the path from plus to minus through $\beta$,
\begin{equation}\label{eq:preal}
 p=-2\int_a^\ell\frac{b(m+y)}{\sqrt{y^2-a^2}}\dd y,
 \qquad p_0=-2\arcosh(\ell/a).
\end{equation}
The negative sign records travel toward the branch on the plus sheet;
the return on the minus sheet contributes the same integral.

\begin{theorem}[Complete transverse and boundary bounds]\label{thm:payment}
Put $C_p=4LR+2KR^2$. With the real contributions of
Proposition~\ref{prop:cutallocation},
\begin{align}
 |p-cp_0|&\le C_p,\label{eq:pbound}\\
 |E_{\cut}+\pi|c|^2p_0|&\le\pi|c|C_p,\label{eq:cutbound}\\
 |E_{\phys}-|c|^2(\Eo+\pi p_0)|&\le C_E+\pi|c|C_p.
 \label{eq:physbound}
\end{align}
These bounds retain the full physical boundary, not only a selected
bounded subtotal.
\end{theorem}
\begin{proof}
By \eqref{eq:preal} and $|b(m+y)-b_0|\le Ly$,
\[
 |p-b_0p_0|\le2L\int_a^\ell\frac{y}{\sqrt{y^2-a^2}}\dd y
 =2L\sqrt{\ell^2-a^2}\le4LR.
\]
Since $\arcosh x\le x$ for $x\ge1$,
\[
 |b_0-c||p_0|\le\frac K2a^2\arcosh(\ell/a)
 \le\frac K2\ell^2\le2KR^2.
\]
This proves \eqref{eq:pbound}. Now $q=2\pi\ii c$ gives
$E_{\cut}=\im(\bar q p)/2=-\pi\re(\bar c p)$.
Subtract the corresponding model to prove \eqref{eq:cutbound}.
Finally $E_{\phys}=\E-E_{\cut}$; Theorem~\ref{thm:remainder} and
\eqref{eq:cutbound} prove \eqref{eq:physbound}.
\end{proof}

Both model boundary payments still have logarithmic growth. The bounded
remainders locate that growth precisely and keep the algebra complete;
they do not make the endpoint energy finite. In particular, boundedness
of $G$ in an unquantified family is neither assumed nor needed in the
proof of the estimates.

\section{A complete moment-controlled heat-event chart}\label{sec:heat}

The preceding estimates apply to a holomorphic pair chart. We now derive
all of its constants from the complete canonical source at an assumed
ordinary double. This construction is conditional local analysis: it
does not assert the existence of such an event for the theta field.

\subsection{Complete-source analyticity and derivative bounds}
Put $\lambda_n=\pi n^2$. The $n$th summand in \eqref{eq:Phi} is
\[
 \Phi_n(u)=\lambda_n e^{5u}(2\lambda_n e^{4u}-3)
                        e^{-\lambda_n e^{4u}}>0
 \qquad(u\geq0),
\]
because $2\lambda_n e^{4u}-3\geq2\pi-3>0$.
For the positive leading terms the successive ratio satisfies
\[
 \frac{(n+1)^4e^{-\pi(n+1)^2e^{4u}}}
      {n^4e^{-\pi n^2e^{4u}}}
 \leq16e^{-3\pi}=:q_0<1.
\]
Dropping the negative part only for this upper bound and summing the
geometric majorant gives
\begin{equation}\label{eq:sourceglobal}
 0<\Phi(u)\leq\frac{2\pi^2}{1-q_0}
                          e^{9u-\pi e^{4u}}.
\end{equation}
Consequently, for every real $S$, $Y\geq0$, and nonnegative integer $j$,
\begin{equation}\label{eq:generalmoment}
 M_{j,Y}(S):=\int_0^\infty u^j e^{Su^2+Yu}\Phi(u)\dd u
                                                        <\infty.
\end{equation}
The decay of $e^{-\pi e^{4u}}$ dominates every displayed quadratic,
linear, and polynomial factor. This also gives a common integrable
majorant on every compact subset of the two complex variables $(t,z)$.
Thus \eqref{eq:Hdef} defines a jointly entire function, its derivatives
can be taken under the complete integral, and
\begin{equation}\label{eq:momentderivative}
 |\partial_z^jH(t,z)|\leq M_{j,Y}(S)
 \quad\text{if }\re t\leq S,\quad |\im z|\leq Y.
\end{equation}
Indeed the derivatives of cosine are cosine or sine with the factor
$u^j$, and both trigonometric functions have modulus at most
$e^{|\im z|u}$. Time differentiation adds two powers of $u$ and
gives \eqref{eq:heat}. The source is real, so $H(t,z)$ is real for
real $t,z$. These statements use all modes and the whole half-line.
An explicit finite acquisition ledger for the moments is given in
Appendix~\ref{sec:tails}.

\subsection{The finite ordinary-double chart}
Assume
\begin{equation}\label{eq:event}
 T\geq0,\quad x_*\in\R,\quad
 H(T,x_*)=H_z(T,x_*)=0,\qquad
 A=\tfrac12H_{zz}(T,x_*)\in\R\setminus\{0\}.
\end{equation}
Use the shifted field and complete moments
\begin{equation}\label{eq:shiftedfield}
 F_\tau(x)=H(T+\tau,x_*+x),\quad h(x)=F_0(x),\quad
 M_j=M_{j,1}(T+1)\quad(j\geq0).
\end{equation}
On the complex bidisk $|\tau|\leq1$, $|x|\leq1$, formula
\eqref{eq:momentderivative} bounds the $j$th spatial derivative of
$F_\tau$ by $M_j$. In particular the following constants are finite
and strictly positive:
\begin{align}
 r&=\min\{1,|A|/(4M_3)\},\label{eq:radius}\\
 \sigma&=\min\{1,3|A|r^2/(8M_2),|A|r/(4M_3),|A|/(4M_4)\}.
 \label{eq:timewindow}
\end{align}
They depend on the event and its curvature. No lower bound uniform in
height, time, or the event is asserted.

\begin{theorem}[Explicit finite event chart]\label{thm:eventchart}
For every complex $|\tau|\leq\sigma$, the field $F_\tau$ has exactly
two zeros in $|x|<r$, counted with multiplicity, and obeys the outer floor
\begin{equation}\label{eq:outerfloor}
 |F_\tau(x)|\geq m_{\rm out}:=3|A|r^2/8
                                      \qquad(|x|=r).
\end{equation}
There is exactly one critical point $k(\tau)$ in this disk. It depends
holomorphically on $\tau$, and the two zeros are simple for
$0<|\tau|\leq\sigma$. For real $0<\tau\leq\sigma$ they are real
and can be ordered as $\alpha<k(\tau)<\beta$. With
\begin{equation}\label{eq:eventmidpoint}
 m=\frac{\alpha+\beta}{2},\qquad a=\frac{\beta-\alpha}{2}>0,
\end{equation}
one has
\begin{equation}\label{eq:splitbounds}
 \frac{6\tau}{5}\leq a^2\leq\frac{10\tau}{3},\qquad
 |m|\leq\frac{28M_3}{27|A|}\tau.
\end{equation}
In particular the finite physical cover is a smooth bordered annulus
for each such positive time, with both physical boundaries retained.
\end{theorem}
\begin{proof}
Taylor's formula with integral remainder gives $h(x)=x^2g(x)$ and
\begin{equation}\label{eq:gcontrol}
 |g(x)-A|\leq M_3|x|/6\leq|A|/24
                                         \qquad(|x|\leq r).
\end{equation}
Thus $g$ is nonvanishing on the closed disk. The only zero of $h$
there is its double zero at the origin, and on the circle
$|h|\geq3|A|r^2/4$. Integrating
$\partial_\tau F_\tau=-F_\tau''$ along the complex time segment gives
\begin{equation}\label{eq:timefieldchange}
 |F_\tau(x)-h(x)|\leq M_2|\tau|\leq3|A|r^2/8.
\end{equation}
The outer floor follows by subtraction. Since the right side is
strictly less than the boundary lower bound for $|h|$, Rouch\'e's
theorem gives the count two.

On the entire closed disk and time window, spatial and time integration
of the derivative bounds gives
\begin{equation}\label{eq:curvaturechart}
 |F_\tau''(x)-2A|\leq M_3r+M_4\sigma\leq|A|/2.
\end{equation}
On $|x|=r$, Taylor's theorem in space and
$\partial_\tau F_\tau'=-F_\tau'''$ likewise give
\[
 |F_\tau'(x)-2Ax|\leq M_3r^2/2+M_3\sigma
                       \leq3|A|r/8<2|A|r.
\]
Rouch\'e's theorem applied to $F_\tau'$ gives exactly one critical
point, counted with multiplicity. Its second derivative is nonzero by
\eqref{eq:curvaturechart}. The local implicit functions therefore
patch, by uniqueness, to one holomorphic function $k(\tau)$ on the
time disk, with $k(0)=0$.

Set $v(\tau)=F_\tau(k(\tau))$. The chain rule, the critical-point
equation, and the heat equation imply
\begin{equation}\label{eq:criticalvalue}
 v'=-F_\tau''(k(\tau)),\quad v(0)=0,\quad
 |v(\tau)+2A\tau|\leq|A||\tau|/2.
\end{equation}
The last inequality follows by integrating
\eqref{eq:curvaturechart} along the complex segment. It implies
$|v(\tau)|\geq3|A||\tau|/2>0$ when $\tau\ne0$. A multiple zero
would have to be the unique critical point, and hence would force
$v=0$. Thus both zeros are simple throughout the punctured time disk.

Now let $\tau$ be positive and real. Reality and uniqueness give a real
critical point. On the real spatial interval the function $f=F_\tau/A$
satisfies
\begin{equation}\label{eq:realconvexity}
 3/2\leq f''\leq5/2,\qquad
 -5\tau/2\leq f(k)\leq-3\tau/2.
\end{equation}
The required endpoint signs follow from more than the modulus floor:
\eqref{eq:gcontrol} gives $h(\pm r)/A\geq3r^2/4$, and
\eqref{eq:timefieldchange} gives
$|F_\tau(\pm r)-h(\pm r)|/|A|\leq3r^2/8$.
Hence $f(\pm r)\geq3r^2/8>0$. Strict convexity now yields one
real zero on each side of $k$, and these exhaust the count two.

For either root $\gamma$, Taylor's theorem about $k$ gives
\[
 \tfrac34|\gamma-k|^2\leq-f(k)
                       \leq\tfrac54|\gamma-k|^2.
\]
Combining this with \eqref{eq:realconvexity}, each of the distances
$k-\alpha$ and $\beta-k$ lies between $\sqrt{6\tau/5}$ and
$\sqrt{10\tau/3}$. Their arithmetic mean is $a$, proving its
two-sided estimate in \eqref{eq:splitbounds}.

To control the midpoint, first use
$|F_\tau'(0)|\leq M_3\tau$, which follows from the time derivative
of $F_\tau'$ and the event equation. The real curvature floor then
gives
\[
 |k|\leq\frac{2M_3}{3|A|}\tau.
\]
Expand $F_\tau(m+a)=F_\tau(m-a)=0$ through second order at $m$.
The two cubic remainders have magnitudes at most $M_3a^3/6$ each.
Subtracting the expansions gives $|F_\tau'(m)|\leq M_3a^2/6$.
Another application of the curvature floor yields
\[
 |m-k|\leq\frac{M_3a^2}{9|A|}
       \leq\frac{10M_3}{27|A|}\tau.
\]
Adding the last two estimates proves the midpoint bound. Finally the
outer floor and simple zero count give the regular finite cover, and
Lemma~\ref{lem:cylinder} identifies it as the asserted annulus.
\end{proof}

\subsection{Deflation and derivative constants from the same source}
For the remainder of this section $0<\tau\leq\sigma$ is real and
$F_\tau(x)=(x-\alpha)(x-\beta)V(x)$.

\begin{lemma}[Explicit deflated-unit estimates]\label{lem:eventunit}
On the complex disk $|x|\leq r$, set $\nu=3|A|/4$. Then
\begin{align}
 |V-A|&\leq|A|/4,\qquad |V|\geq\nu,\label{eq:unitfloor}\\
 |V'|&\leq M_3/6,\qquad |V''|\leq M_4/12.\label{eq:unitderivatives}
\end{align}
For a coherent holomorphic choice $b=V^{-1/2}$ the constants
\begin{equation}\label{eq:eventLK}
 L=\frac{M_3}{12\nu^{3/2}},\qquad
 K=\frac{M_3^2}{48\nu^{5/2}}+\frac{M_4}{24\nu^{3/2}}
\end{equation}
satisfy $|b|\leq\nu^{-1/2}$, $|b'|\leq L$, and $|b''|\leq K$.
At the root midpoint one has the sharper estimate
\begin{equation}\label{eq:Jmidpoint}
 |V(m)-A|\leq J\tau,
 \qquad J=\frac{14M_3^2}{27|A|}+\frac{23M_4}{36}.
\end{equation}
\end{lemma}
\begin{proof}
Let $\Delta=\{(s,t):s,t\geq0,\ s+t\leq1\}$. Twice applying the
fundamental theorem of calculus along complex line segments gives
\begin{equation}\label{eq:simplexunit}
 V(x)=\int_\Delta
 F_\tau''\bigl(s\alpha+t\beta+(1-s-t)x\bigr)\dd s\dd t.
\end{equation}
To verify the normalization, for distinct $x,\alpha,\beta$ first
integrate in $t$ and then in $s$. The result is
\[
 \frac1{\beta-x}
 \left\{\frac{F_\tau(\alpha)-F_\tau(\beta)}{\alpha-\beta}
       -\frac{F_\tau(\alpha)-F_\tau(x)}{\alpha-x}\right\}
 =\frac{F_\tau(x)}{(x-\alpha)(x-\beta)}.
\]
The cases $x=\alpha,\beta$ follow by removable continuation.
The disk is convex, so every evaluation point in
\eqref{eq:simplexunit} belongs to the same disk on which our source
derivative bounds hold. The elementary simplex integrals are
\begin{align*}
 \int_\Delta1&=\tfrac12,&
 \int_\Delta(1-s-t)&=\tfrac16,&
 \int_\Delta(1-s-t)^2&=\tfrac1{12},\\
 \int_\Delta(t-s)&=0,&
 \int_\Delta(t-s)^2&=\tfrac1{12}.&&
\end{align*}
For example these follow by successive integration of monomials, or
from $\int_\Delta s^it^j(1-s-t)^k\dd s\dd t
=i!j!k!/(i+j+k+2)!$. Integrating
\eqref{eq:curvaturechart} proves \eqref{eq:unitfloor}; differentiating
\eqref{eq:simplexunit} once and twice in $x$ and using the next two
simplex integrals proves \eqref{eq:unitderivatives}.

The unit has no zero, so it has a holomorphic square root on the disk.
The derivative identities
\[
 b'=-\tfrac12V^{-3/2}V',\qquad
 b''=\tfrac34V^{-5/2}(V')^2-\tfrac12V^{-3/2}V''
\]
give exactly \eqref{eq:eventLK}. On the real interval, $V$ is real
and \eqref{eq:unitfloor} makes its sign the sign of $A$. Therefore
$b$ has a fixed nonzero phase there. In particular its arcsine average,
which determines the core period in Section~\ref{sec:remainders},
is nonzero.

At $x=m$ the simplex evaluation point is $m+a(t-s)$. Expand
$F_\tau''$ at $m$. Its linear term integrates to zero, and the
quadratic remainder and the last simplex integral give
\[
 |V(m)-\tfrac12F_\tau''(m)|\leq M_4a^2/24.
\]
Time integration and spatial integration also give
$|F_\tau''(m)-2A|\leq M_4\tau+M_3|m|$. Combining with
\eqref{eq:splitbounds} yields
\begin{align*}
 |V(m)-A|
 &\leq\tfrac12M_4\tau+\tfrac12M_3|m|+M_4a^2/24\\
 &\leq\left(\frac{14M_3^2}{27|A|}
            +\frac12M_4+\frac5{36}M_4\right)\tau,
\end{align*}
which is \eqref{eq:Jmidpoint}.
\end{proof}

The deflated-unit floor \eqref{eq:unitfloor} is a conclusion of the
event equations and complete derivative bounds. It is not a lower floor
for $F_\tau$, which vanishes at each branch. Theorem~\ref{thm:remainder}
therefore applies on this chart with the explicitly derived constants
$\nu,L,K,r$; its constants \eqref{eq:Cosc} and \eqref{eq:CE} do not
depend on the shrinking pair separation.

\subsection{A finite coarse upper bound}
\begin{proposition}\label{prop:coarseheat}
For $0<\tau\leq\sigma$ the complete energy obeys
\begin{equation}\label{eq:Bcoarse}
 \E_\tau\leq B(\tau):=
 \frac8{3|A|}\left[4\pi+2\pi\log\left(
                   2r\sqrt{\frac5{6\tau}}\right)\right].
\end{equation}
In particular $[0,B(\sigma/2)]$ encloses all energies on the bounded
window $\sigma/2\leq\tau\leq\sigma$.
\end{proposition}
\begin{proof}
Set $d=\beta-\alpha=2a$. Split the base disk into its intersections
with $B(\alpha,d/2)$ and $B(\beta,d/2)$ and the complement. The two
small disks are disjoint except at their tangent point. On the first,
$|x-\beta|\geq d/2$, so its contribution to the integral of
$1/(|x-\alpha||x-\beta|)$ is bounded by the integral of
$2/(d|x-\alpha|)$ over the full radius-$d/2$ disk, which is $2\pi$.
The same bound holds at the other root. On the complement the two
distances $p,q$ lie between $d/2$ and $2r$. Use
$1/(pq)\leq(p^{-2}+q^{-2})/2$ and bound each resulting integral by
its full centered annulus integral. This gives
\[
 \int_{|x|<r}\frac{\dd A_x}{|x-\alpha||x-\beta|}
                  \leq4\pi+2\pi\log(4r/d).
\]
The two-sheet energy has the additional factor two, and
$|V|\geq3|A|/4$. Finally $d\geq2\sqrt{6\tau/5}$ by
\eqref{eq:splitbounds}. These facts give \eqref{eq:Bcoarse}.
The right side decreases with $\tau$, proving the window statement.
\end{proof}

\subsection{A strictly sharper shared upper certificate}
Write the explicit constant \eqref{eq:Cosc} on the present chart as
\begin{equation}\label{eq:heatCosc}
 C_{\osc}=8\pi rL/\sqrt\nu+8\pi r^2L^2.
\end{equation}
Define
\begin{align}
 D&=2r\sqrt{5/6},\qquad c_*=\log2+\tfrac14,\label{eq:sharpD}\\
 K_*&=\max\{0,3|A|C_{\osc}/(2\pi)+7c_*-16\},\label{eq:sharpK}\\
 s_*&=\min\{\sigma,D^2,\nu/(6J),D^2e^{-2(K_*+1)}\}>0.
 \label{eq:sharpwindow}
\end{align}
The canonical source has $J>0$, since its moments are positive. If the
same argument is used for a polynomial calibration with $J=0$, the
division term in \eqref{eq:sharpwindow} is omitted.

\begin{theorem}[Strict improvement of the coarse certificate]
\label{thm:sharpheat}
For $0<\tau\leq s_*$ a valid complete upper bound is
\begin{equation}\label{eq:Usharp}
 \E_\tau\leq U(\tau):=
 4\pi\left(\frac1{|A|}+\frac{J\tau}{|A|\nu}\right)
               \arsinh(D/\sqrt\tau)+C_{\osc}.
\end{equation}
It satisfies the explicit strict comparison
\begin{equation}\label{eq:strictBU}
 B(\tau)-U(\tau)\geq\frac{2\pi}{3|A|}>0.
\end{equation}
Thus $[0,U(\tau)]$ is strictly contained in $[0,B(\tau)]$ on this
declared window.
\end{theorem}
\begin{proof}
The root midpoint is real and the two roots are interior, so
$r-|m|>a$ and $r+|m|\leq2r$. The upper shifted-disk bound
\eqref{eq:shiftedsandwich} gives
\[
 \Eo\leq4\pi\arsinh((r+|m|)/a)
       \leq4\pi\arsinh(D/\sqrt\tau),
\]
where the last step uses \eqref{eq:splitbounds}. Also
\eqref{eq:unitfloor} and \eqref{eq:Jmidpoint} imply
\[
 \frac1{|V(m)|}-\frac1{|A|}
 =\frac{|A|-|V(m)|}{|A||V(m)|}
 \leq\frac{J\tau}{|A|\nu}.
\]
Apply the $b(m)$ version of Theorem~\ref{thm:remainder}, with
remainder \eqref{eq:Cosc}, to obtain \eqref{eq:Usharp}.

For $x\geq1$ the elementary inequalities
$\sqrt{1+y}\leq1+y/2$ and $\log(1+y/4)\leq y/4$,
with $y=x^{-2}\leq1$, give
\begin{equation}\label{eq:arsinhestimate}
 \arsinh x=\log x+\log(1+\sqrt{1+x^{-2}})
                          \leq\log x+\log2+\tfrac14.
\end{equation}
Let $\ell=\log(D/\sqrt\tau)$. From \eqref{eq:sharpwindow},
$\ell\geq K_*+1$, $D/\sqrt\tau\geq1$, and
$1+J\tau/\nu\leq7/6$. The coarse certificate is
$B=(2\pi/(3|A|))(16+8\ell)$, whereas
\eqref{eq:Usharp} and \eqref{eq:arsinhestimate} give
\[
 U\leq\frac{2\pi}{3|A|}(7\ell+7c_*)+C_{\osc}.
\]
Subtracting gives
\[
 B-U\geq\frac{2\pi}{3|A|}
       \left[\ell+16-7c_*-\frac{3|A|C_{\osc}}{2\pi}\right]
       \geq\frac{2\pi}{3|A|},
\]
by the definition of $K_*$.
\end{proof}

This is a strict improvement over a specified coarse certificate. Both
bounds are available to direct integration supplied with the same roots,
unit, and source bounds. Thus the theorem gives a sharper shared direct
baseline; it supplies no ordering of all boundary and direct methods.

\subsection{The resulting degeneration is still logarithmic}
\begin{corollary}[Ordinary-double energy scale]\label{cor:heatasymptotic}
For the fixed event chart above,
\begin{equation}\label{eq:heatasymptotic}
 \E_\tau=\frac{2\pi}{|A|}\log(1/\tau)+O(1)
                                    \qquad(\tau\downarrow0).
\end{equation}
The bounded remainder may depend on the event and its chart constants.
\end{corollary}
\begin{proof}
By \eqref{eq:splitbounds}, $a\longrightarrow0$,
$\log(1/a)=\tfrac12\log(1/\tau)+O(1)$, and $m=O(\tau)$.
In particular, for sufficiently small time both $r-|m|$ and $r+|m|$
lie in a fixed positive interval about $r$. Apply the lower and upper
shifted-disk bounds \eqref{eq:shiftedsandwich} and the elementary
large-argument asymptotics used in \eqref{eq:quadasymp}. They give
\[
 \Eo=4\pi\log(1/a)+O(1)
      =2\pi\log(1/\tau)+O(1).
\]
The unit estimate \eqref{eq:Jmidpoint} and the floor $|V(m)|\geq\nu$
give the two-sided estimate
\[
 \left|\frac1{|V(m)|}-\frac1{|A|}\right|
                       \leq\frac{J\tau}{|A|\nu}.
\]
Finally Theorem~\ref{thm:remainder}, using \eqref{eq:Cosc}, gives
$\E_\tau=|V(m)|^{-1}\Eo+O(1)$. The normalization error is
$O(\tau\log(1/\tau))$, hence bounded and tending to zero. This proves
\eqref{eq:heatasymptotic}.
\end{proof}

The positive logarithmic coefficient is consistent with the exact
quadratic control. Neither \eqref{eq:Bcoarse} nor \eqref{eq:Usharp}
remains bounded as $\tau$ tends to the singular fiber. A source estimate
that excludes the event would have to add information not supplied by
this generic ordinary-double chart.

\section{Analytic consequences for finite enclosures}\label{sec:enclosures}

The preceding statements are exact. A numerical acquisition must also
control the uncertainty of every term that it uses. This section records
the analytic bounds behind such an acquisition; it does not prescribe
an implementation or assume that the boundary expression is cheaper.

\subsection{Primitive products and period products}
Suppose the cut polygon is decomposed into finitely many oriented arcs
with real parameters $s$. On each arc write $\omega=B(s)\dd s$ and let
$F$ denote the coherent primitive in Proposition~\ref{prop:stokes}.
Suppose approximations satisfy
\[
 |F-\widehat F|\le\delta_F(s),\qquad
 |B-\widehat B|\le\delta_B(s).
\]
An independent quadrature enclosure for
$(\ii/2)\sum\int\widehat F\overline{\widehat B}\dd s$ has error
$E_{\rm quad}$. Then its real part $\widehat E$ satisfies
\begin{equation}\label{eq:producterror}
 |\E-\widehat E|\le E_{\rm quad}+
 \frac12\sum\int\left(\delta_F|\widehat B|
       +|\widehat F|\delta_B+\delta_F\delta_B\right)\dd s.
\end{equation}
Indeed expand $(\widehat F+\Delta_F)
\overline{(\widehat B+\Delta_B)}$ and apply the triangle inequality.
The primitive error includes the basepoint and accumulated integral
errors; it cannot be reset independently on different boundary arcs.
Paired-edge cancellation is usable only with the corresponding common
primitive and jump enclosure.

Likewise, if $|q-\widehat q|\le\epsilon_q$ and
$|p-\widehat p|\le\epsilon_p$, then
\begin{equation}\label{eq:perioderror}
 \left|\im(\bar q p)-\im(\overline{\widehat q}\widehat p)\right|
 \le|\widehat q|\epsilon_p+|\widehat p|\epsilon_q
                         +\epsilon_q\epsilon_p.
\end{equation}
This follows by the same expansion. In a collision chart $|p|$ grows
logarithmically, so a fixed absolute uncertainty in $q$ is amplified.
A bounded exact remainder in Theorem~\ref{thm:payment} alone does not
imply a bounded acquisition error.

\subsection{Simple-root neighborhoods}
If $H(z)=(z-\alpha)U(z)$ on $|z-\alpha|\le\eta$, with
$0<\nu_1\le|U|\le M_1$, then direct polar integration gives
\begin{equation}\label{eq:cap}
 \frac{4\pi\eta}{M_1}
 \le\E(H;D(\alpha,\eta))
 \le\frac{4\pi\eta}{\nu_1}.
\end{equation}
Thus the neighborhood of a simple zero has an explicit finite payment.
Equivalently the coordinate \eqref{eq:branchregular} removes the local
singularity. An exterior floor on $H$ is never used as an interior floor
across a root. When two roots are present, the full coordinate formula
\eqref{eq:directU} is a direct integration route using exactly the same
unit and root information as the period formulas.

\subsection{A complete analytic approximation and exact deflation}
Let $H$ be holomorphic on $|z-C|\le\mathcal R$ with $|H|\le M$ there,
and let $0<R<\mathcal R$, $\rho=R/\mathcal R$. For the degree-$N$
Taylor polynomial $P_N$ about $C$, Cauchy's coefficient estimate gives,
when $N\ge2$,
\begin{align}
 \epsilon_0&:=\sup_{|z-C|\le R}|H-P_N|
 \le M\frac{\rho^{N+1}}{1-\rho},\label{eq:tay0}\\
 \epsilon_1&:=\sup_{|z-C|\le R}|H'-P_N'|
 \le\frac M{\mathcal R}\rho^N
       \left(\frac{N+1}{1-\rho}+\frac\rho{(1-\rho)^2}\right),
 \label{eq:tay1}\\
 \epsilon_2&:=\sup_{|z-C|\le R}|H''-P_N''|
 \le\frac M{\mathcal R^2}\rho^{N-1}
 \left(\frac{N(N+1)}{1-\rho}
 +\frac{(2N+1)\rho}{(1-\rho)^2}
 +\frac{\rho(1+\rho)}{(1-\rho)^3}\right).
 \label{eq:tay2}
\end{align}
For example the last formula is the sum of
$(N+1+k)(N+k)\rho^k$, using
$\sum\rho^k=(1-\rho)^{-1}$,
$\sum k\rho^k=\rho(1-\rho)^{-2}$, and
$\sum k^2\rho^k=\rho(1+\rho)(1-\rho)^{-3}$.
The first two are the corresponding zeroth and first derivatives of the
geometric tail. For the complete theta field, $M$ can be supplied by a
complete source moment; \eqref{eq:xi} is also an exact alternative at
$t=0$.

Let $\alpha,\beta$ be the exact simple roots, enclosed independently,
and define the polynomial divided difference
$Q_N(z)=P_N[\alpha,\beta,z]$. Even though $P_N$ need not vanish at the
exact roots, the complete function does vanish at them, and its divided
difference is $V(z)=H[\alpha,\beta,z]$. The elementary simplex identity proved in
Section~\ref{sec:heat} therefore yields
\begin{equation}\label{eq:deflationerror}
 |V(z)-Q_N(z)|\le\frac{\epsilon_2}{2}
 \quad(z,\alpha,\beta\in\overline{D(C,R)}).
\end{equation}
Indeed the second derivative of $H-P_N$ is integrated over a simplex of
area $1/2$; the disk is convex. This argument does not divide a field
error by $|(z-\alpha)(z-\beta)|$ near a root. It retains the exact
removable continuation and a uniform error bound. Root enclosures must
still be propagated through $Q_N$ and the coordinate map; replacing the
roots by their midpoints would not justify \eqref{eq:deflationerror}.

All these bounds apply before selecting a mesh or an arithmetic package.
They specify the error information required by either comparison route.
In particular, taking the intersection of two proved intervals is a
shared improvement available to both routes.

\section{Recorded numerical support}\label{sec:support}

The finite analytic framework has been applied to one actual regular
theta fiber, without an RH assumption. The recorded outward-ball
calculation \cite{ThetaRecord}, using Arb arithmetic \cite{Arb}, treats
$t=0$ and $|z-35|\le8$. A complete boundary winding count certifies
exactly two interior roots, each separately enclosed as simple and real:
\[
 \alpha\simeq28.26945028346938758091450396712,\qquad
 \beta\simeq42.04407927754310998525695918780.
\]
The certificate uses the exact normalization \eqref{eq:xi}, with complete
analytic Taylor and divided-difference tails, rather than a finite source
head treated as the full field. Its native boundary interval is
\begin{equation}\label{eq:numBoundary}
 [28770927.163097,\ 33752291.350598],
\end{equation}
and its native direct-coordinate interval is
\begin{equation}\label{eq:numDirect}
 [28728270.232214,\ 34655897.029090].
\end{equation}
The endpoints are outward rational bounds as printed. Both intervals
meet the recorded common relative-diameter target, and the first lies
strictly inside the second. The result supports a finite acquisition
on a genuine regular fiber. It compares the two specified native
implementations; a direct method permitted to reuse all the same
identities and observations inherits their intersection. Neither the
interval containment nor the recorded timing difference establishes an
intrinsic advantage for the energy method. Detailed arithmetic,
accounting, and replay instructions remain in \cite{ThetaRecord}.
No theorem above uses these numerical intervals.

\section{What has been achieved and what remains open}\label{sec:open}

\subsection{The present analytic conclusion}
The project now has a complete finite-domain identity, a full quadratic
calibration, explicit bounds for the energy and both boundary payments
after quadratic subtraction, and a source-controlled ordinary-event
chart. The source tail estimates in Appendix~\ref{sec:tails} make the
constants effective without replacing the complete canonical source.
The sharper upper bound in Section~\ref{sec:heat} improves a specified
coarse bound on an explicit event-dependent interval. These statements
hold with all spatial boundaries, both sheets, the mixed term, and the
relative period retained.

The same analysis explains its current limitation. At an ordinary event,
\[
 \E_\tau=\frac{2\pi}{|A|}\log(1/\tau)+O(1),
\]
and the derived upper bounds have this permitted growth. Their constants
are finite on each fixed chart but are not a theta-specific obstruction
to that chart. The polynomial heat family $z^2-2\tau$ satisfies the local
mechanism exactly. Consequently ordinary holomorphicity, the heat
equation, real symmetry, and a controlled nonvanishing unit cannot alone
supply the desired exclusion.

\subsection{Endpoint control and coverage}
For a collision exclusion through reciprocal energy, one would need a
new estimate for the unchanged complete source, for example
\begin{equation}\label{eq:missingbound}
 \sup_{0<\tau\le\tau_0}\E(H(T+\tau,\cdot);D)<\infty
\end{equation}
on a fixed disk containing the collision, with hypotheses not already
assuming the conclusion. The logarithmic lower bound would then give
a contradiction. No estimate of the form \eqref{eq:missingbound} is
proved here. Alternatively, a source inequality that forces a sufficiently
strong upper control of the relative period, or of the full exterior
payment after the common constant is fixed, could address the same gap.
It must be checked against the complete quadratic calculation.

Even such an ordinary-event estimate would need a coverage argument.
The present theorem assumes two real simple roots in a fixed disk and
starts from a hypothetical real nondegenerate double event. It does not
treat every height, higher multiplicities, arbitrary nonreal pairs,
or a moving spatial boundary. An individual simple nonreal zero still
has finite reciprocal energy by Proposition~\ref{prop:multi}; regular
fiber finiteness by itself does not locate zeros on the real axis.
Any implication toward RH or toward RH plus simplicity therefore needs
an additional global bridge, with its reality regime, endpoint passage,
and uniform constants stated separately. The usual canonical heat-flow
context is described in \cite{RT}; it does not provide the missing
energy estimate merely by a change of language.

\subsection{Controls on a proposed source-specific advance}
A future source-specific inequality should state exactly which canonical
properties it uses and test their separate roles. A positive finite
weighted-square or finite-mode control may preserve useful positivity
while failing the complete normalized reflection identity and the
unchanged canonical coefficients. A reflected primitive control may
preserve evenness or a transform identity while losing the canonical
square-index sequence. An Euler-product modification may change the
arithmetic source itself. Such controls answer different questions;
none should be said to preserve the entire canonical class merely
because one identity survives.

The relevant next analytic task is thus to isolate a complete-source
input that controls degeneration beyond the generic chart bounds.
The finite identities and paid remainders here specify where that input
must enter. A compact-surface or Jensen-based reserve approach would
likewise require its own quantitative passage from finite bordered
surfaces; the compact bilinear relation alone does not remove the
physical boundary or supply the needed uniformity.

\appendix
\section{Canonical normalization and source reflection}\label{sec:normalization}

For completeness we derive the classical normalization used in
\eqref{eq:xi}. Let
\[
 \psi(x)=\sum_{n=1}^\infty e^{-\pi n^2x},\qquad
 \Theta(x)=1+2\psi(x)\quad(x>0).
\]
The Fourier transform of $e^{-\pi x v^2}$, with kernel
$e^{-2\pi\ii v\xi}$, is $x^{-1/2}e^{-\pi\xi^2/x}$.
Poisson summation for this rapidly decaying Gaussian therefore gives
\begin{equation}\label{eq:Jacobitheta}
 \Theta(x)=x^{-1/2}\Theta(1/x),\qquad
 \psi(x)=x^{-1/2}\psi(1/x)+\tfrac12(x^{-1/2}-1).
\end{equation}
Set $f(u)=e^u\psi(e^{4u})$ for real $u$. Direct differentiation of the
locally uniformly convergent series gives
\begin{equation}\label{eq:Phioperator}
 \frac{f''(u)-f(u)}8
 =2e^{9u}\psi''(e^{4u})+3e^{5u}\psi'(e^{4u})=\Phi(u).
\end{equation}
The theta identity gives $f(-u)=f(u)+\sinh u$. Applying
$\partial_u^2-1$, which annihilates $\sinh u$, proves
$\Phi(-u)=\Phi(u)$. Differentiating at zero also gives
$-f'(0)=f'(0)+1$, or $f'(0)=-1/2$.
These identities concern the complete series.

For $\re s>1$, absolute convergence and the gamma integral give
\[
 \pi^{-s/2}\Gamma(s/2)\zeta(s)
 =\int_0^\infty\psi(x)x^{s/2-1}\dd x.
\]
Split at $x=1$, apply \eqref{eq:Jacobitheta} on $(0,1)$, and substitute
$x\mapsto1/x$ in its decaying part. The elementary remaining integral is
$1/(s-1)-1/s=1/(s(s-1))$. Hence
\begin{equation}\label{eq:completedMellin}
 \pi^{-s/2}\Gamma(s/2)\zeta(s)
 =\frac1{s(s-1)}+
 \int_1^\infty\psi(x)
 \left(x^{s/2-1}+x^{(1-s)/2-1}\right)\dd x.
\end{equation}
The last integral is entire in $s$ by its exponential decay. Multiplying
by $s(s-1)/2$ thus continues $\xi$ as an entire function. Put
$v=2s-1$ and $x=e^{4u}$ in \eqref{eq:completedMellin}. This gives
\begin{equation}\label{eq:ximellin}
 \xi(s)=\frac12+(v^2-1)\int_0^\infty f(u)\cosh(vu)\dd u.
\end{equation}
Twice integrating by parts, using the super-exponential decay of $f,f'$
at $+\infty$ and $\sinh(0)=0$, yields
\begin{align*}
 \int_0^\infty(f''-f)\cosh(vu)\dd u
 &=-f'(0)+(v^2-1)\int_0^\infty f(u)\cosh(vu)\dd u\\
 &=\xi(s).
\end{align*}
All these integrals converge for every complex $v$, so analytic
continuation introduces no new restriction. Finally take
$s=(1+\ii z)/2$, for which $v=\ii z$, and use
\eqref{eq:Phioperator}. The result is
\[
 H(0,z)=\int_0^\infty\Phi(u)\cos(zu)\dd u
       =\frac18\xi\bigl((1+\ii z)/2\bigr),
\]
as claimed. This derivation also fixes the factor $1/8$ used by the
energy and numerical records; a scalar change of $H$ would rescale its
reciprocal energy by the inverse scalar modulus.

\section{Effective complete-source moment tails}\label{sec:tails}

The moments used in the event chart retain the entire theta source and
the entire integration half-line. This appendix gives explicit upper
bounds for replacing them by finitely many modes on a finite interval.
For $S\ge0$, $Y\ge0$, and an integer $j\ge0$, recall the notation
\begin{equation}\label{eq:tails-moment}
 M_{j,Y}(S)=\int_0^\infty u^j e^{Su^2+Yu}\Phi(u)\dd u.
\end{equation}
Write $\lambda_n=\pi n^2$ and split the source in
\eqref{eq:Phi} into its complete individual summands,
\begin{equation}\label{eq:tails-summand}
 \Phi_n(u)=\lambda_n e^{5u}(2\lambda_ne^{4u}-3)
                         e^{-\lambda_ne^{4u}}.
\end{equation}
For $u\ge0$, the factor $2\lambda_ne^{4u}-3$ is at least
$2\pi-3>0$. Thus every $\Phi_n$ is positive, and
\begin{equation}\label{eq:tails-leading}
 0<\Phi_n(u)\le 2\lambda_n^2e^{9u-\lambda_ne^{4u}}.
\end{equation}
The smaller term in \eqref{eq:tails-summand} is discarded only in
this upper estimate; it remains part of the finite source below.

Choose an integer $N\ge0$ and a real $U\ge1$. Put $k=N+1$ and define
the exact finite head
\begin{equation}\label{eq:tails-head}
 M^{N,U}_{j,Y}(S)=\sum_{n=1}^{N}
       \int_0^U u^j e^{Su^2+Yu}\Phi_n(u)\dd u.
\end{equation}
An empty sum is understood as zero. Positivity and Tonelli's theorem,
initially allowing infinite values, give the disjoint decomposition
\begin{align}
 M_{j,Y}(S)=M^{N,U}_{j,Y}(S)
 &+\sum_{n=k}^\infty\int_0^\infty
                  u^j e^{Su^2+Yu}\Phi_n(u)\dd u\notag\\
 &+\sum_{n=1}^{N}\int_U^\infty
                  u^j e^{Su^2+Yu}\Phi_n(u)\dd u.
 \label{eq:tails-decomposition}
\end{align}
The first omitted piece comprises all omitted modes on the whole
half-line. The second comprises only retained modes on their exterior
rays. Their index--integration regions do not overlap.

\begin{proposition}[All omitted modes]\label{prop:tails-modes}
Set
\begin{equation}\label{eq:tails-mode-parameters}
 \rho_k=\left(\frac{k+1}{k}\right)^4e^{-\pi(2k+1)},
 \qquad b_k=4\lambda_k-9-Y.
\end{equation}
If $b_k>0$ and $S\le8\lambda_k$, then
\begin{equation}\label{eq:tails-modes}
 \sum_{n=k}^\infty\int_0^\infty
       u^j e^{Su^2+Yu}\Phi_n(u)\dd u
 \le L_{j,Y}(S;N):=
 \frac{2\lambda_k^2e^{-\lambda_k}j!}
                  {(1-\rho_k)b_k^{j+1}}.
\end{equation}
In particular $\rho_k<1$ for every $k\ge1$.
\end{proposition}

\begin{proof}
For $n\ge k$ and $u\ge0$, the ratio of consecutive leading terms in
\eqref{eq:tails-leading} is
\[
 \frac{\lambda_{n+1}^2e^{-\lambda_{n+1}e^{4u}}}
      {\lambda_n^2e^{-\lambda_ne^{4u}}}
 =\left(\frac{n+1}{n}\right)^4
                 e^{-\pi(2n+1)e^{4u}}
 \le\rho_k.
\]
Both the rational factor and the exponential factor are bounded above
by their values at $n=k$, $u=0$. Also
$\rho_k\le16e^{-3\pi}<1$. Geometric summation therefore gives
\begin{equation}\label{eq:tails-geometric}
 \sum_{n=k}^\infty\Phi_n(u)
 \le\frac{2\lambda_k^2}{1-\rho_k}
                 e^{9u-\lambda_ke^{4u}}.
\end{equation}
For $u\ge0$, Taylor's inequality for the positive exponential series is
\begin{equation}\label{eq:tails-exp-inequality}
 e^{4u}\ge1+4u+8u^2.
\end{equation}
Consequently
\begin{align*}
 Su^2+(9+Y)u-\lambda_ke^{4u}
 &\le-\lambda_k-(4\lambda_k-9-Y)u
                         +(S-8\lambda_k)u^2\\
 &\le-\lambda_k-b_ku.
\end{align*}
Multiply \eqref{eq:tails-geometric} by $u^je^{Su^2+Yu}$ and
integrate. The elementary integral
\[
 \int_0^\infty u^je^{-bu}\dd u=\frac{j!}{b^{j+1}}
 \qquad(b>0)
\]
follows by scaling and $j$ integrations by parts. It gives exactly
\eqref{eq:tails-modes}.
\end{proof}

\begin{proposition}[Retained exterior rays]\label{prop:tails-rays}
For each retained index $1\le n\le N$, suppose that
\begin{equation}\label{eq:tails-ray-parameters}
 b_n(U)=4\lambda_ne^{4U}-9-2SU-Y>0,
 \qquad S\le8\lambda_ne^{4U}.
\end{equation}
Then
\begin{align}
 \int_U^\infty u^je^{Su^2+Yu}\Phi_n(u)\dd u
 &\le R_{n,j,Y}(S;U),\label{eq:tails-ray-bound}\\
 R_{n,j,Y}(S;U)
 &:={2\lambda_n^2}
 e^{SU^2+(9+Y)U-\lambda_ne^{4U}}
 \sum_{h=0}^{j}\binom{j}{h}
       \frac{U^{j-h}h!}{b_n(U)^{h+1}}.
 \label{eq:tails-ray}
\end{align}
\end{proposition}

\begin{proof}
Set $u=U+v$, where $v\ge0$. Applying
\eqref{eq:tails-exp-inequality} to $v$ gives
\begin{align*}
 S(U+v)^2+(9+Y)(U+v)-\lambda_ne^{4U}e^{4v}
 \le{}&SU^2+(9+Y)U-\lambda_ne^{4U}\\
 &-b_n(U)v+(S-8\lambda_ne^{4U})v^2.
\end{align*}
The quadratic coefficient is nonpositive by hypothesis. Thus
\eqref{eq:tails-leading} bounds the exterior integral by
\[
 2\lambda_n^2e^{SU^2+(9+Y)U-\lambda_ne^{4U}}
       \int_0^\infty(U+v)^je^{-b_n(U)v}\dd v.
\]
The binomial theorem and the same elementary exponential integrals give
\[
 \int_0^\infty(U+v)^je^{-bv}\dd v
 =\sum_{h=0}^{j}\binom{j}{h}
                  \frac{U^{j-h}h!}{b^{h+1}},
\]
which proves \eqref{eq:tails-ray-bound}--\eqref{eq:tails-ray}.
\end{proof}

\begin{corollary}[Complete finite enclosure]\label{cor:tails-complete}
Under the hypotheses of Propositions~\ref{prop:tails-modes} and
\ref{prop:tails-rays},
\begin{equation}\label{eq:tails-complete}
 M^{N,U}_{j,Y}(S)\le M_{j,Y}(S)
 \le M^{N,U}_{j,Y}(S)+L_{j,Y}(S;N)
                       +\sum_{n=1}^{N}R_{n,j,Y}(S;U).
\end{equation}
For each fixed $S,Y,j$, admissible $N,U$ exist and the sum of the
displayed tail bounds can be made arbitrarily small.
\end{corollary}

\begin{proof}
Insert the two upper bounds into the positive disjoint identity
\eqref{eq:tails-decomposition}. To verify the final assertion, first
increase $N$. Then $b_k=4\pi k^2-9-Y$ is positive,
$8\pi k^2\ge S$, $\rho_k\to0$, and the right-hand side of
\eqref{eq:tails-modes} tends to zero because of $e^{-\pi k^2}$.
With such an $N$ fixed, increase $U$. For the finitely many retained
indices, the term $e^{4U}$ dominates both $U$ and the constant $S$;
hence \eqref{eq:tails-ray-parameters} eventually holds simultaneously.
In \eqref{eq:tails-ray}, the negative exponent
$-\lambda_ne^{4U}$ dominates $SU^2+(9+Y)U$, while the finite sum has
only polynomial powers of $U$ and inverse positive powers of $b_n(U)$.
Thus every $R_{n,j,Y}(S;U)$ tends to zero. This proves the assertion
and also proves finiteness of \eqref{eq:tails-moment} directly.
\end{proof}

These bounds estimate the complete positive moments. They do not impose
separate reflection identities on a source head or tail, and they make
no replacement of the canonical coefficient sequence. In particular, the
finite head in \eqref{eq:tails-head} retains both terms of every
summand in \eqref{eq:tails-summand}.

\section*{Acknowledgements and draft status}
This manuscript was drafted for Edward Baker with substantial assistance
from OpenAI's Codex, recorded model \texttt{gpt-6.1-sol}, reasoning effort
\emph{ultra}. LLM assistance was used in derivation, checking, exposition,
and preparation of the accompanying computations. Separate agents of the
same model supplied internal checks; these are not independent specialist
or human validation. The manuscript consolidates the local research notes
and retains their stated scope. Classical identities are attributed below;
no mathematical priority claim is made. The numerical support has a
replayable outward-ball record, while the theorems and proofs in the main
text do not depend on a numerical sign or timing result.
\begin{thebibliography}{9}
\bibitem{RT}
B.~Rodgers and T.~Tao,
\emph{The de Bruijn--Newman constant is non-negative},
Forum of Mathematics, Pi \textbf{8} (2020), e6, 62 pp.
\href{https://doi.org/10.1017/fmp.2020.6}{doi:10.1017/fmp.2020.6}.
The normalization used here is Introduction, equations (1)--(4).

\bibitem{Grava}
T.~Grava, \emph{Riemann Surfaces and Integrable Systems},
lecture notes, 14 February 2021, Section 3.1.2.
\url{https://people.sissa.it/~grava/RS2021.pdf}.

\bibitem{Arb}
F.~Johansson, \emph{Arb: Efficient Arbitrary-Precision Midpoint-Radius
Interval Arithmetic}, 2016.
\url{https://arxiv.org/abs/1611.02831}.

\bibitem{Notes}
E.~Baker research project, \emph{Reciprocal Hodge energy},
Notes 01--03, 22 October 2026. Prepared for Edward Baker with substantial
LLM assistance, \texttt{gpt-6.1-sol}, effort ultra.
Accompanying \repo\ working repository, directory
\path{papers/investigations-4/09_reciprocal_hodge_energy/notes/}.
The proofs are developed independently of these records in the present text.

\bibitem{ThetaRecord}
E.~Baker research project, \emph{Actual regular theta fiber comparison},
Note 04, certificate \path{03_ACTUAL_THETA_COMPARISON_20261022.json},
and Review 03, 22 October 2026. Prepared for Edward Baker with substantial
LLM assistance, \texttt{gpt-6.1-sol}, effort ultra.
Accompanying repository directory
\path{papers/investigations-4/09_reciprocal_hodge_energy/}.
The certificate and its three replay scripts are in \path{numerics/};
the implementation-specific comparison and all error budgets are documented
in Note 04 and the fixed comparison contract.
\end{thebibliography}
\end{document}
